% What the index indicates, limitations and scope decisions
\section{What the index indicates}\label{sec:meaning}

An index is only useful if a high value means something one can act on. A high \Fd\ marks a crop and place
where three things coincide: the growing area is already under heat or water stress, the crop loses yield
quickly under that stress, and people depend on the crop for nutrients. Read this way, our results say
the following.

\begin{enumerate}
\item \textbf{The most vulnerable crop is not the most heat-sensitive one.} Soybean has the steepest
literature heat slope, yet it ranks low because in all seven of its countries the monthly flowering-season
temperature stays below its threshold (zero heat exposure) and is eaten
mainly as oil. Wheat is only moderately sensitive, but it is grown in the hottest and driest plains of the
panel and is the largest single food source of zinc and iron in most panel countries. Vulnerability is about
\emph{overlap}, which neither a yield study nor a nutrient study shows on its own.
\item \textbf{The nutrient at stake is iron.} Iron is the only nutrient whose national supply falls below
the requirement in any panel country (\nIronShort\ of \nAdqCountries; energy, protein, fat and zinc in none), and the crops that supply most iron (wheat; pulses in South Asia and East Africa) are among
those whose climate stress is highest or rising fastest. Climate risk therefore falls on a nutrient that is
already short, not on calories.
\item \textbf{Concentrated diets amplify the signal.} Where a few crops supply most nutrients (India gets
\divIndIronTopShare\% of its iron from wheat; its supply of energy is spread over an effective
\divIndKcal\ food items compared with \divWorldKcal\ for the world as a whole), a loss in one vulnerable
crop removes a larger share of the nutrient.
\item \textbf{Irrigation is the hidden buffer.} Rice scores low largely because most of its area is
irrigated, which our water measure counts as unstressed. The index therefore assumes that irrigation water
keeps flowing; groundwater depletion or reservoir shortfalls would move rice up the ranking, and our
monthly data cannot see the hot days that matter for rice.
\item \textbf{The future changes magnitudes more than the order.} Under one model and one scenario the
ranking of crops is largely the same in 2041--2060, but vulnerability rises for wheat (through heat and
water in similar measure) and common bean (mostly through heat during flowering). Pulses, a key protein and iron source for vegetarian diets, are the
group whose score rises fastest relative to its present level.
\end{enumerate}

\paragraph{What it points to.} Read as a screening tool, the index points to where adaptation would protect
most nutrition: for wheat in India, whose present score is mainly water but whose projected increase comes
almost entirely from heat during flowering (\texttt{future\_components.csv}), heat tolerance and sowing dates
as well as water; water management for wheat in Russia and the USA, where both today's score and most of the
projected increase are water-driven; drought management for rainfed maize in southern Africa; heat-tolerant beans in
tropical Brazil and East Africa; and, on the demand side, diversifying iron sources so that one crop's bad
year does not decide iron supply. These are directions suggested by where scores are high, not tested
interventions.

\section{Limitations and scope decisions}\label{sec:limits}
\begin{itemize}
\item \textbf{Monthly climate data} hide hot days, so heat exposure is a lower bound (rice and soybean
score near zero on heat for this reason). \textbf{Water stress} is a screening ratio of rainfall to crop water
need, with no soil storage; this overstates stress for crops grown on stored soil moisture (chickpea in
India, which sets the water reference) and treats irrigated area as unstressed.
\item \textbf{New-crop parameters} are fits to chamber or field experiments, not values stated by their
sources; the rapeseed heat slope and lentil $K_y$ are proxies; the bean threshold comes from one
heat-susceptible Andean variety and is applied to FAO ``beans, dry'', which in India and Myanmar include
tropical \emph{Vigna} pulses. These are tested (Table~\ref{tab:robust}) and wheat's first place survives.
\item \textbf{Crop maps are from 2010} (SPAM) and \textbf{calendars} are partly gap-filled; chickpea, lentil
and rapeseed use a documented wheat-based rule. Pigeon pea (tur), a major pulse in South India, is not
among the nine crops.
\item \textbf{One climate model and one scenario} (MRI-ESM2-0, SSP2-4.5): the future results show
direction under that model, not the range of possible futures.
\item \textbf{Nutrition} uses national supply, not individual intake; one composition value per item; and
recommended intakes that overstate shortfall when compared with mean supply. The future yield step is an
accounting estimate without CO$_2$ fertilisation, adaptation or market response, and the empirical
regressions do not confirm it pair by pair.
\item \textbf{Scope decisions.} The future yield-to-nutrition step was implemented as a conditional
screening estimate rather than a forecast. The CO$_2$ effect on nutrient content is reported only as an
additional pathway \citep{myers2014}, not built into the index. Diet diversity is measured as the diversity
of national food \emph{supply} (FAO) and, for Hyderabad, of household consumption (MoSPI), not individual
dietary diversity scores.
\end{itemize}

\section{Conclusion}
Across nine crops, the Climate--Nutrition Vulnerability Index identifies \topEq\ as the crop whose climate
stress most overlaps with its nutritional role, a result that holds across \robTopFfour\ of \robNFfour\
robustness rankings and is stable to 2041--2060 under one climate model. Cereals as a group are most
vulnerable, pulses are rising fastest, and the nutrient most at stake is iron. The index is most useful as a
map of where climate stress and dietary dependence coincide, which is where adaptation and diet
diversification would protect nutrition most.
